\documentclass[11pt]{article}
\usepackage{amsmath,amssymb,amsthm,mathtools,hyperref}

\begin{document}
\section*{Universal microbial growth laws}
\label{prob:056}
\label{rank:220}
\par\smallskip

\noindent\textbf{Domain.} Comparative microbial physiology \quad \textbf{Retained status/source context.} Open domain-of-validity and mechanistic universality question; established growth laws have restricted regimes.
\paragraph{Definitions and assumptions.} For taxon $a$ and condition $E$, measure balanced exponential growth rate $\mu_{a,E}$, ribosome mass fraction $\phi_{R,a,E}$, effective peptide elongation rate $k_{a,E}$ and mean cell volume $V_{a,E}$. Let $\ell_a$ be independently measured characteristic protein length and $V_{0,a}$ a specified reference volume. Normalize $z=\mu\ell_a/k$ and $v=V/V_{0,a}$. Restrict initially to viable cultures without nutrient exhaustion; distinguish nutrient limitation, translation inhibition and stress conditions.
\paragraph{Problem statement.} Determine the maximal experimentally supported class of taxa and conditions for which a common relation holds,
\[\phi_{R,a,E}=g(z_{a,E};\vartheta_a)+r_{a,E},\qquad |r_{a,E}|\le\varepsilon,\]
with a small predeclared taxon-specific parameter vector $\vartheta_a$ rather than a separately fitted curve for every condition. Derive $g$ from allocation or synthesis constraints and prospectively predict where departures exceed tolerance. Ask separately whether $v$ obeys an additional relation conditioned on replication and division physiology.
\paragraph{Additional binding assumptions and resolution criterion.} Choose a prespecified nested family of taxon/condition classes or report all inclusion-maximal supported classes; there need not be a unique largest one. Units of elongation rate and protein length must make $z$ dimensionless. Parameter count, fitting data and extrapolation rules are fixed before validation. Finite-panel success does not establish a universal relation outside that panel.

\paragraph{Alternative formulations.} (i) Test a linear ribosomal-fraction law within one bacterium under nutrient modulation, using an intercept interpreted as inactive ribosomes. (ii) Seek dimensionless collapse across bacterial and archaeal taxa after accounting for active ribosome fraction, maintenance and temperature. These are different strength claims; no equivalence with a cell-size law is assumed.
\paragraph{Obstructions and partial progress.} Empirical relations and coarse-grained mechanistic explanations exist in E. coli, so finding any ribosome-growth correlation is not an open problem. Translation inhibitors can generate different slopes from nutrient changes, and volume depends on cell-cycle regulation in addition to protein synthesis. A convincing universality claim must specify its exclusions and predict failures rather than discard them retrospectively. Cross-taxon assays also need comparable mass-fraction and activity definitions; a species-dependent measurement bias can create or erase an apparent data collapse.
\paragraph{Sources (supplied provenance; not all rechecked).}
\begin{itemize}
\item Georgouli et al., Multi-scale models of whole cells: progress and challenges (2023). \url{https://pmc.ncbi.nlm.nih.gov/articles/PMC10661945/} (Supplied review; broad background, with abstract and metadata corroborated through PubMed.)
\item Scott et al., Interdependence of cell growth and gene expression: origins and consequences (2010). \url{https://doi.org/10.1126/science.1192588} (Growth laws demonstrated for specified bacterial conditions.)
\item Scott et al., Emergence of robust growth laws from optimal regulation of ribosome synthesis (2014). \url{https://doi.org/10.15252/msb.20145379} (Mechanistic growth-law analysis in E. coli.)
\item Georgouli et al., Multi-scale models of whole cells: progress and challenges (2023). \url{https://pubmed.ncbi.nlm.nih.gov/38020904/} (Broad modelling limitations; no theorem of universal prediction.)
\end{itemize}


\section*{Binding catalogue conventions}
This edition revises both supplied catalogues and preserves original problem numbers 1--330. It adds one globally unique label \verb|\label{rank:x}| per problem, where smaller $x$ means greater provisional difficulty. A rank label is an anchor at that problem's section: \verb|\ref{rank:x}| returns the section number, not the literal difficulty number. The visible \texttt{rank:x} text and the index give the difficulty number itself.

The entries are \emph{parameterized mathematical research schemes}, including empirical prediction, causal identification, historical reconstruction, theory and engineering. They are not 330 certified unsolved mathematical conjectures. A rigorous instance fixes the objects and restrictions below; changing those objects changes the question. Established identities, finite calculations, non-identification constructions and solved special cases are not being posed anew as open. Some residual tasks concern informative bounds, biological realization or external validation rather than a new theorem.

Every main definition and problem statement received a mathematical editorial review. Each entry now has an additional binding-conditions paragraph; direct corrections are separately recorded in the accompanying audit. The alternative formulations inherit the applicable conventions, but are different targets unless equivalence is proved. Where an old broad phrase conflicts with an explicit restriction below it, the explicit restriction controls.

\textbf{Literature boundary.} Thirteen selective primary-literature checks cover eighteen entries in this revision; their access levels and scope are recorded explicitly. The supplied bibliography and its earlier source-access descriptions are retained as provenance, not represented as newly verified. Neither this review nor a source's description proves the absence of later solutions. The date of this edition is 13 October 2026; there is no claim of exhaustive current-status certification for all 330 entries.

\subsection*{Binding conventions for a concrete instance}
\paragraph{Objects, quantifiers and non-vacuity.}
An instance declares measurable state, input, output and parameter spaces; a fixed candidate family; admissible initial laws, interventions and environments; a finite horizon or specified asymptotic regime; observation and sampling laws; and performance metrics, positive normalization scales and positive tolerances. These are chosen before fitting or testing, not adjusted after seeing failures. Shared parameters must genuinely be shared across the declared validation domain. Finite dimension alone does not restrict arbitrary real-number encodings: specify a finite library, bounded description length, or regularity/complexity bounds. Exclude constant or no-signal solutions through a fixed informative-output or baseline-improvement requirement where appropriate. A universally quantified result cannot be established by choosing one trivial favorable instance.

\paragraph{Probability laws and data.}
Every expectation or probability refers to the specified stochastic law, including initial uncertainty, process noise, observation noise and sampling as appropriate. A deterministic equation is not a probability model without a declared random input or initial law. ``Observed'' targets in population-level formulas mean population observables $g_j^*$, not noisy estimates asserted to be exact. Let $\widetilde g_j$ be independent validation estimates and suppose
\[
\mathbb{P}\{\lvert\widetilde g_j-g_j^*\rvert\le r_j\text{ for every }j\}\ge1-\alpha.
\]
Then an uncertainty-aware sufficient certification rule for positive target tolerances $\varepsilon_j$ is
\[
\max_j\frac{|\widehat g_j-\widetilde g_j|+r_j}{\varepsilon_j}\le1.
\]
On that simultaneous confidence event it implies every population error is at most its tolerance, by the triangle inequality. Vector/norm versions use corresponding valid simultaneous bounds. Training uncertainty must also be included when the claim averages over repeated training data. Evaluation on finitely many conditions does not establish a continuous-domain supremum without covering and regularity bounds.

\paragraph{Dynamics and well-posedness.}
For ODEs and SDEs give coefficient domains, initial laws and existence/uniqueness conditions in the intended solution sense, or keep the entire admitted solution set. For stochastic jump models specify nonnegative off-diagonal rates, diagonal normalization, rate orientation and non-explosion. For PDEs specify spatial domains, boundary/initial conditions, conservation or source terms, and a weak or classical solution class. Moving-domain concentration equations include dilution/transport terms. A stationary law is used only for a time-homogeneous augmented model under existence assumptions; a stationary approximation to a driven system must be justified. Absorbing extinction states are retained. Survival-conditioned, quasi-stationary and unconditional laws are distinct and must not be substituted for one another.

\paragraph{Identification and causal meaning.}
Two parameter values are observationally equivalent when they induce the same full distribution under all allowed observation/intervention designs, not merely the same likelihood value for one sample. A target is identified only if constant on the equivalence class. Gauge symmetries, label switching and reference conventions are quotiented or fixed. Finite-data compatible sets use a declared confidence, likelihood or discrepancy rule; an empty set signals model rejection. Bayesian credible sets are not automatically confidence sets. A prior or regularizer does not create data identification. Bounds intended to resolve a question must be informative at a predeclared width or decision threshold.

A causal contrast uses well-defined intervention versions and a stated causal model. Random assignment, or explicit consistency, overlap and exchangeability assumptions, is required for causal interpretation from observations. Interference and time-varying confounding are represented when present. Fixing a mediator by intervention differs from conditioning on its observed value. Marginal potential-outcome laws do not determine individual cross-world effects or an arbitrary joint coupling. Ethically admissible designs, rather than unconstrained interventions, define the experimental menu.

\paragraph{Metrics, units, optimization and limits.}
Specify the sampling or reference measure for densities and integrals, the log base, and dimensionless arguments of logarithms. Total variation is $\sup_A|P(A)-Q(A)|$; Wasserstein distances require the declared ground metric and appropriate moment conditions; KL divergence uses its stated absolute-continuity convention. All normed readouts are scaled consistently, and error probabilities use $0<\alpha<1$. A ratio has a positive denominator or an explicit undefined/zero convention. Thermodynamic entropy production is nonnegative in expectation under the relevant hypotheses; individual stochastic realizations need not be. An infimum is not an attained minimum without compactness/coercivity and continuity or another existence argument. Fix the order of large-population, long-time and vanishing-mutation limits. Stability claims declare perturbation space and symmetries.

\subsection*{How the provisional difficulty order was assigned}
Four editorial components are scored from 1 (lower burden) to 5 (higher burden): (a) formal or inferential difficulty, (b) access to evidence and identification burden, (c) mechanistic/computational integration, and (d) independent validation and transfer. Their unweighted total orders the entries. For selected entries in the highest bands, an explicit editorial priority breaks equal totals; all other ties are broken by original problem number. Equal totals are \emph{not} evidence of reliably different true difficulties. Scores are judgments, not measurements, completion-time estimates or consensus rankings.

The object ranked is an informative, nondegenerate resolution of the \emph{scoped} scientific target, with the stated budget and validation domain. A sharp useful impossibility result can qualify; an already known identity, vacuous bound, arbitrary favorable model or lookup table cannot. Rank 330 is merely the least difficult relative assessment in this catalogue, not a claim of an easy biological solution. Broad historical titles do not automatically outrank precisely formulated mechanistic tasks.



\end{document}
